\chapter{Cohomological invariants and depth} \label{III}

\section{Recollections} \label{III.1}

We shall state some definitions and results which the reader will find
for example in Chapter~I of the course given by J.-P\ptbl Serre at the
Coll\`ege de France in 1957--58.\nde{\label{noteserre} the re-edition of
Serre's text (Serre~J.-P., \emph{Alg\`ebre locale. {M}ultiplicit\'es},
{Cours au Coll\`ege de France, 1957--1958, written by Pierre Gabriel,
second edition, Lect. Notes in Math.}, vol.~11, Springer-Verlag, 1965)
no longer contains the proofs of these statements. One may refer to
(Bourbaki~N., \textit{Alg\`ebre commutative}, Masson), as Serre himself
suggests.}

\begin{definition} \label{def:III.1.1}
Let \(A\) be a ring, (commutative with unit element as in all that
follows) and let \(M\) be an \(A\)-module (unitary as in all that
follows), one calls:
\begin{itemize}
\item
annihilator of \(M\) and one denotes by \(\Ann M\) the set of
\(a\in A\) such that for every \(m\in M\) one has \(am=0\).

\item
support of \(M\) and one denotes by \(\Supp M\) the set of prime
ideals \(\pp \) of \(A\) such that the localized module \(M_\pp\) is
nonzero.

\item
``assassin of \(M\)'' or ``set of prime ideals associated with \(M\)''
and one denotes by \(\Ass M\) the set of prime ideals \(\pp \) of \(A\)
such that there exists a nonzero element of~\(M\) whose annihilator is
\(\pp \).
\end{itemize}

If \(\ideal a\) is an ideal of \(A\), we shall denote by
\(\rr(\ideal a)\) the radical of \(\ideal a\) in \(A\), \ie the set of
elements of \(A\) a power of which lies in \(\ideal a\).
\end{definition}

The following results are valid if one assumes that
\emph{\(A\) is noetherian and \(M\) of finite type}.

\setcounter{subsection}{0}

\begin{proposition} \label{III.1.1}
\begin{enumeratei}
\item
\(\Ass M\) is a finite set.

\item
For an element of \(A\) to annihilate a nonzero element of \(M\), it is
necessary and sufficient that it belong to one of the ideals associated
with \(M\).

\item
The radical of the annihilator of \(M\), \(\rr(\Ann M)\), is the
intersection of the associated ideals of \(M\) which are minimal (for
the inclusion relation in \(\Ass M\)).
\end{enumeratei}
\end{proposition}

\begin{proposition} \label{III.1.2}
Let \(\pp \) be a prime ideal of \(A\); the following assertions are
equivalent:
\begin{enumeratei}
\item
\(\pp \in \Supp M\).

\item
There exists \(\qq \in \Ass M\) such that \(\qq \subset\pp \).

\item
\(\pp \supset \Ann M\).

\item[\textup{(iii bis)}]
\(\pp \supset \rr(\Ann M)\).
\end{enumeratei}
\end{proposition}

\begin{proposition} \label{III.1.3}
Let \(N\) be an \(A\)-module of finite type; one has the formula:
\[
\Ass \Hom_A(N, M) = \Supp N \cap \Ass M.
\]
\end{proposition}

\section{Depth} \label{III.2}

Throughout this section, \(A\) denotes a commutative ring, \(I\) an
ideal of \(A\), \(M\)~and~\(N\) two \(A\)-modules. One will denote by
\(X\) the prime spectrum of \(A\) (one will not use its structure sheaf
in this section) and by \(Y\) the variety of \(I\),
\(Y = \Supp(A/I) = \{\pp \in\nobreak X, \ \pp \supset I\}\).

\begin{lemma} \label{III.2.1}
Suppose that \(A\) is noetherian and that the modules \(M\) and \(N\)
are of finite type. Suppose moreover that \(\Supp N=Y\). Then the
following assertions are equivalent:
\begin{enumeratei}
\item
\(\Hom_A(N, M)=0\).
\item
\(\Supp N \cap \Ass M = \emptyset\).
\item
The ideal \(I\) is not a zero-divisor in \(M\), which means that for
every \(m\in M\), \(Im=0\) implies \(m=0\).
\item
There exists in \(I\) an \(M\)-regular element. (An element \(a\) of
\(A\) is said to be \(M\)-regular if the homothety of ratio \(a\) in
\(M\) is injective.)

\item
For every \(\pp \in Y\), the maximal ideal \(\pp A_\pp\) of the local
ring \(A_\pp\) is not associated with \(M_\pp\). In a formula:
\(\pp A_\pp \notin \Ass M_\pp\).
\end{enumeratei}
\end{lemma}

\begin{proof}
(i) \SSI (ii) since \(\Ass \Hom_A(N, M)=\emptyset\) is equivalent to
(ii) by proposition \Ref{III.1.3} and to (i) by an easy consequence of
proposition \Ref{III.1.2}.

(iii) \ALORS (ii) by contradiction: ``there exists
\(\pp \in \Supp N \cap \Ass M\)'' implies that \(\pp \supset I\) and
that there exists \(m\in M\) whose annihilator is \(\pp \), hence
\(Im \subset \pp m=0\), which contradicts (iii).

(iv) \ALORS (iii) trivially.

(ii) \SSI (iv) since \(\Supp N = Y\), hence (ii) means that \(I\) is
not contained in any ideal associated with \(M\), or equivalently
(since the ideals associated with \(M\) are prime and finite in
number), that \(I\) is not contained in the union of the ideals
associated with \(M\). Now, by proposition \Ref{III.1.1}~(ii), this set
is the set of elements of \(A\) which are not \(M\)-regular.

(i) \ALORS (v); indeed, if \(\Hom_A(N, M)=0\) and if \(\pp \in Y\), one
deduces, by virtue of the formula
\[
\bigl(\Hom_A(N, M) \bigr)_\pp = \Hom_{A_\pp} \bigl(N_\pp, M_\pp \bigr),
\]
that \(\Hom_{A_\pp} \bigl(N_\pp, M_\pp \bigr)=0\), hence, thanks to
proposition \Ref{III.1.3},
\[
\Supp N_\pp \cap \Ass M_\pp = \emptyset,
\]
now \(\pp A_\pp \in \Supp N_\pp, \) hence \( \pp A_\pp \notin \Ass M_\pp\).

(v) \ALORS (i); indeed, if \(\pp \in \Ass M\), there exists \(m\in M\)
whose annihilator is \(\pp\), hence the canonical image of \(m\) in
\(M_\pp\) is nonzero, hence its annihilator is an ideal which contains
\(\pp\), hence \(\pp A_\pp\), hence is equal to it. The ideal
\(\pp A_\pp\) is therefore associated with \(M_\pp\), hence
\(\pp \notin Y\) by (v), whence (i).
\end{proof}

We shall work with these conditions, replacing the
functor \(\Hom\) by its derived functors.

\begin{theorem} \label{III.2.2}
Let \(A\) be a commutative ring, \(I\) an ideal of \(A\), \(M\) an
\(A\)-module. Let \(n\) be an integer.

\begin{enumeratea}
\item
If there exists a sequence \(f_1, \ldots, f_{n+1}\) of elements of
\(I\) which forms an \(M\)-regular sequence (\ie if \(f_1\) is
\(M\)-regular and if \(f_{i+1}\) is regular in
\(M/(f_1, \ldots, f_i)M\) for \(i\leq n\)), then for every
\(A\)-module \(N\) annihilated by a power of \(I\), one has:
\[
\Ext_A^i(N, M)=0 \text{ for } i\leq n.
\]

\item
If moreover \(A\) is noetherian, if \(M\) is of finite type, and if
there exists an \(A\)-module \(N\) of finite type such that
\(\Supp N = \variete(I)\) and such that \(\Ext_A^i(N, M)=0\) for
\(i\leq n\), then there exists a sequence \(f_1, \ldots, f_{n+1}\) of
elements of \(I\) which is \(M\)-regular.
\end{enumeratea}
\end{theorem}

Let us first prove a), by induction. If \(n<0\) the statement is empty.

If \(n\geq 0\), suppose that a) has been proved for \(n'<n\); by
hypothesis there exists \(f_1\in I\) which is \(M\)-regular. Let us
denote by \(f_1^i\) multiplication by \(f_1\) in \(\Ext_A^i(N, M)\) and
by \(f_1^M\) multiplication by \(f_1\) in \(M\). The sequence
\begin{equation} \label{eq:III.2.1}
0 \to M \lto{f_1^M} M \to M/f_1M \to 0
\end{equation}
is exact, hence so is the sequence:
\[
\Ext_A^{i-1} (N, M/f_1M) \lto{\delta}\Ext_A^i (N, M) \lto{f_1^i} \Ext_A^i (N, M).
\]
By hypothesis \(I^rN=0\), hence \(f_1^0\) and nilpotent; \(\Ext^i\) is
a universal functor, and the same is true of \(f_1^i\) for every \(i\).
Moreover, there exists a regular sequence in \(M/f_1 M\) which has \(n\)
elements, hence, by the induction hypothesis,
\[
Ext_A^{i-1}(N, M)=0 \text{ if } i\leq n-1.
\]
One deduces that if \(i\leq n\), \(f_1^i\) is both nilpotent and
injective, hence \(\Ext_A^i(N, M)=0\).

Let us prove b), likewise by induction. If \(n<0\), the statement is
empty.

If \(n=0\), b) follows from the assertion (i) \ALORS (iv) of lemma
\Ref{III.2.1}.

If \(n>0\), by b) for \(n=0\), there exists an element \(f_1\in I\)
which is \(M\)-regular; from the exact sequence (\Ref{eq:III.2.1}), one
deduces the exact sequence:
\begin{equation} \label{eq:III.2.2} \Ext_A^{i-1} (N, M) \to \Ext_A^{i-1} (N, M/f_1M) \to \Ext_A^i (N, M).
\end{equation}
One concludes that the hypotheses of b) are satisfied for the
module \(M/f_1M\) and for the integer \(n-1\). By the induction
hypothesis, there exists a sequence of \(n\) elements of \(I\) which is
regular for \(M/f_1M\), which implies that there exists a
sequence of \(n+1\) elements of~\(I\), beginning with \(f_1\), and which
is \(M\)-regular.

This theorem invites us to generalize in the following way the classical
definition of the depth of a module of finite type over a noetherian
ring:

\begin{definition} \label{III.2.3}
Let \(A\) be a commutative ring with unit element, let \(M\) be an
\(A\)-module, let \(I\) be an ideal of \(A\). One calls \(I\)-depth of
\(M\), and one denotes by \(\prof_I M\), the supremum in
\(\NN\cup\{+\infty\}\) of the set of natural numbers \(n\) which are
such that for every \(A\)-module of finite type \(N\) annihilated by a
power of \(I\), one has
\[
\Ext_A^i(N, M)=0 \text{ for every } i< n.
\]
\end{definition}

One deduces from the preceding theorem that if \(n\) is the supremum of
the lengths of the \(M\)-regular sequences of elements of \(I\), one
has \(n\leq\prof_I M\).

\noindent More precisely:

\begin{proposition} \label{III.2.4}
Let \(A\) be a commutative ring, \(I\) an ideal of \(A\) and let \(M\)
be an \(A\)-module, and finally let \(n\in \NN\). Consider the
assertions:
\begin{enumerate}
\item
\(n\leq \prof_I M\).
\item
For every \(A\)-module of finite type \(N\) which is annihilated by a
power of \(I\), one has:
\[
\Ext_A^i(N, M)=0 \text{ for } i<n.
\]
\item
There exists an \(A\)-module of finite type \(N\) such that
\(\Supp N = \variete(I)\) and such that \(\Ext_A^i(N, M)=0\) if
\(i<n\).
\item
There exists an \(M\)-regular sequence of length \(n\) formed of
elements of \(I\).
\end{enumerate}

One has the following logical implications:
\[
\begin{array}{c}\dpl \textup{(1)} \Ssi \textup{(2)} \From (4) \\
\big\Downarrow \\
\textup{(3)} \\
\end{array}
\]
Moreover if \(A\) is noetherian and \(M\) of finite type, these
conditions are equivalent.
\end{proposition}

\begin{proof}
(1) \SSI (2) by definition and (2) \ALORS (3) upon taking \(N=A/I\);
Moreover (4) \ALORS (2) by theorem \Ref{III.2.2} a). Finally, if \(A\)
is noetherian and \(M\) of finite type, (3) \ALORS (4) by theorem
\Ref{III.2.2} b).
\skipqed
\end{proof}

We assume \(A\) noetherian and \(M\) of finite type until the end of
this section.

\begin{corollary} \label{III.2.5}
Let \(f\in I\) be an \(M\)-regular element; one has:
\[
\prof_I M = \prof_I (M/fM) + 1.
\]
\end{corollary}

Indeed, if \(n\leq \prof_I (M/fM)\), there exists a sequence of
elements of \(I\), \(f_1, \ldots, f_n\), which is
\((M/fM)\)-regular; hence the sequence \(f, f_1, \ldots, f_n\) is
\(M\)-regular, hence \(n+1 \leq \prof_I M\), hence
\(\prof_I M\geq \prof_I (M/fM) + 1\). On the other hand, by the exact
sequence (\Ref{eq:III.2.2}), if \(i \leq \prof_I M\), one has
\(\Ext_A^{i-1} (N, M/fM)=0\), hence
\(\prof_I M - 1 \leq \prof_I (M/fM)\).

\begin{corollary} \label{III.2.6}
Every finite \(M\)-regular sequence, formed of elements of \(I\), can
be extended to a maximal \(M\)-regular sequence, whose length is
necessarily equal to the \(I\)-depth of~\(M\).
\end{corollary}

\begin{remark} \label{III.2.7}
One can scarcely restrain oneself from saying that an \(A\)-module is
all the more beautiful the greater its depth. A module whose support
does not meet \(\variete (I)\) is among the most beautiful; indeed, one
can prove that for \(\prof_I M\) to be finite, it is necessary and
sufficient that \(\Supp M \cap \variete(I) \neq \emptyset\).
\end{remark}

\begin{remark} \label{III.2.8}
If \(A\) is a semi-local ring, let \(\rr(A)\) be its radical and
\(k=A/\rr(A)\) its residue ring. The interesting notion of depth is
obtained by taking for \(I\) the radical of~\(A\). We shall therefore
agree to denote simply by \(\prof M\) the \(\rr(A)\)-depth of an
\(A\)-module~\(M\). In this case one recovers the notion of
``homological codimension'', (\cf Serre, \textit{op\ptbl cit.}
note~\eqref{noteserre}, page~\pageref{noteserre}), which was denoted
\(\codimh_A M\), and which is defined as the infimum of the integers
\(i\) such that \(\Ext_A^i(k, M)\neq 0\); indeed
\(\Supp k = \variete(\rr(A))\).
\end{remark}

\begin{proposition} \label{III.2.9}
If \(A\) is noetherian and \(M\) of finite type, one has:
\[
\prof_I M = \inf_{\pp \in \variete (I)} \prof M_\pp.
\]
\end{proposition}

\begin{corollary} \label{III.2.10}
If \(A\) is a noetherian semi-local ring, and if \(M\) is an
\(A\)-module of finite type, one has:
\[
\prof M = \inf_\mm \prof M_\mm,
\]
where \(\mm \) runs through the set of maximal ideals of \(A\).
\end{corollary}

The corollary follows immediately from the proposition \Ref{III.2.9};
indeed the prime ideals which contain the radical are the maximal
ideals.

On the other hand, let \(f\in I\); if \(f\) is \(M\)-regular, if
\(\pp\in X\) and if \(\pp\supset I\), the image \(g\) of \(f\) in
\(A_\pp\) belongs to \(\pp A_\pp\), the maximal ideal of \(A_\pp\);
moreover \(g\) is \(M_\pp\)-regular, as follows from the exact sequence
\begin{equation} \label{eq:III.2.3} 0 \to M_\pp \lto{g'} M_\pp \to (M/fM)_\pp \to 0,
\end{equation}
where \(g'\) denotes the homothety of ratio \(g\) in \(M_\pp\). This
exact sequence also gives that \((M/fM)_\pp\) is isomorphic to
\(M_\pp/gM_\pp\); applying corollary \Ref{III.2.5} to \(M\) and to
\(M_\pp\), one deduces, by induction, that \(\prof_I M\leq \nu(M)\),
where one has set for every \(M\):
\[
\nu(M)=\inf_{\pp\in\variete(I)} \prof M_\pp.
\]
More precisely, still by induction, one knows, if \(f\) is
\(M\)-regular, that \(\nu(M)= \nu(M/fM)+1\); it remains therefore to
prove that if \(\nu(M)\geq 1\), there exists an \(M\)-regular element
in \(I\). Now applying lemma \Ref{III.2.1} to \(M_\pp\), \(A_\pp\) and
\(\pp A_\pp\) for every \(\pp\in \variete(I)\), one sees that
\(\pp A_\pp \notin \Ass M_\pp\), hence, applying lemma \Ref{III.2.1} to
\(A\), \(M\) and \(I\), one has the conclusion.

\begin{proposition} \label{III.2.11}
Let \(u\colon A\to B\) be a homomorphism of noetherian rings. Let~\(I\)
be an ideal of \(A\), \(M\) an \(A\)-module of finite type. Set
\(I_B=I\otimes_A B\) and \(M_B=M\otimes_A B\). If \(B\) is \(A\)-flat,
one has:
\[
\prof_{I_B} M_B \geq \prof_I M;
\]
moreover if \(B\) is faithfully flat over \(A\), one has equality.
\end{proposition}

Indeed, let \(N=A/I\); by flatness one has: \(N\otimes_A B=B/I_B\); set
\(N_B=N\otimes_A B\). Still by flatness and noetherian hypotheses, one
has:
\[
\Ext_B^i (N_B, M_B) = \Ext_A^i(N, M)\otimes_A B,
\]
hence \(\Ext_A^i(N, M) = 0 \) implies \(\Ext_B^i(N_B, M_B)=0\), and the
converse is true if \(B\) is faithfully flat over \(A\).
